\documentclass[aspectratio=169]{beamer}

\usetheme{Madrid}
\usecolortheme{default}

\usepackage{amsmath,amssymb,bm}
\usepackage{physics}
\usepackage{hyperref}
\usepackage{mathtools}

\title{Diffusion Models: Foundations (Chapters 1--6)}
\subtitle{Lecture slides based on arXiv:2510.21890v1 (figures and code omitted)}
\author{}
\date{}

\begin{document}

\frame{\titlepage}

\begin{frame}{Outline}
\tableofcontents
\end{frame}

% =========================================================
\section{Lecture 0: Framing and Notation}
% =========================================================

\begin{frame}{Learning goals}
\begin{itemize}
\item Understand diffusion models as deep generative models that transform noise into data.
\item See three equivalent perspectives:
\begin{itemize}
\item variational (DDPM / hierarchical latent model),
\item score-based (denoising score matching / Score SDE),
\item flow-based (CNF / probability flow ODE / flow matching).
\end{itemize}
\item Derive core training objectives and sampling dynamics.
\item Connect to PDEs: continuity equation and Fokker--Planck.
\end{itemize}
\end{frame}

\begin{frame}{Basic notation}
\begin{itemize}
\item Data: \(x\in\mathbb{R}^D\), data distribution \(p_{\text{data}}(x)\).
\item Model: \(p_\theta(x)\) or sampler \(x = G_\theta(z)\), \(z\sim p_{\text{prior}}\).
\item Time/noise: discrete \(t=0,\dots,T\) or continuous \(t\in[0,1]\).
\item Score: \(s_t(x)\equiv \nabla_x \log p_t(x)\).
\item Expectation \(\mathbb{E}[\cdot]\), KL divergence \(\mathrm{KL}(q\|p)\).
\end{itemize}
\end{frame}

\begin{frame}{Transport viewpoint}
All diffusion-style methods can be seen as learning a time-dependent transformation
\[
p_{\text{prior}} \xrightarrow{\text{learned dynamics}} p_{\text{data}}.
\]
Core object: a \textbf{vector field} (or drift) \(v(x,t)\) that moves samples.
\begin{itemize}
\item DDPM: learns a denoiser / noise predictor.
\item Score methods: learn \(s_t(x)\).
\item Flow matching: learn \(v(x,t)\) directly.
\end{itemize}
\end{frame}

\begin{frame}{How these slides map to Chapters 1--6}
\begin{enumerate}
\item Ch.1: deep generative modeling landscape.
\item Ch.2: diffusion probabilistic models (variational view).
\item Ch.3: score matching and denoising score matching (EBM lineage).
\item Ch.4: continuous-time formulation: SDEs, reverse-time, probability flow ODE.
\item Ch.5: flows, CNFs/Neural ODE, flow matching.
\item Ch.6: unifying principles: conditioning + Fokker--Planck / continuity equations.
\end{enumerate}
\end{frame}

% =========================================================
\section{Chapter 1: Deep Generative Modeling Primer}
% =========================================================

\begin{frame}{Generative modeling problem}
Given i.i.d. samples \(x_i\sim p_{\text{data}}\), learn \(p_\theta\) so that
\[
p_\theta \approx p_{\text{data}}
\]
and we can sample new \(x\sim p_\theta\).
Core tensions:
\begin{itemize}
\item likelihood tractability,
\item sample quality,
\item sampling speed,
\item mode coverage and stability.
\end{itemize}
\end{frame}

\begin{frame}{Explicit vs implicit models}
\begin{itemize}
\item \textbf{Explicit density}: can evaluate (or approximate) \(\log p_\theta(x)\).
\item \textbf{Implicit}: provides sampling procedure but no tractable likelihood.
\end{itemize}
Diffusion models: often trained via likelihood-type bounds or regression objectives, but sampling is iterative/dynamical.
\end{frame}

\begin{frame}{Maximum likelihood}
Maximum likelihood:
\[
\max_\theta \ \mathbb{E}_{p_{\text{data}}}[\log p_\theta(x)].
\]
When direct likelihood is hard:
\begin{itemize}
\item introduce latent variables (VAE),
\item impose invertibility (flow),
\item match scores (score matching),
\item use adversarial objectives (GAN).
\end{itemize}
\end{frame}

\begin{frame}{Divergences and objectives}
Common divergences/metrics:
\begin{itemize}
\item \(\mathrm{KL}(p_{\text{data}}\|p_\theta)\) (maximum likelihood).
\item \(\mathrm{KL}(p_\theta\|p_{\text{data}})\) (mode-seeking behavior).
\item Wasserstein distances and other IPMs (GAN family).
\item Score matching: match \(\nabla_x \log p\) rather than \(p\).
\end{itemize}
Different choices shape sample sharpness vs mode coverage.
\end{frame}

\begin{frame}{Autoregressive (AR) models}
Factorization for \(x=(x_1,\dots,x_D)\):
\[
p_\theta(x)=\prod_{i=1}^D p_\theta(x_i\mid x_{<i}).
\]
Pros: exact likelihood.
Cons: sampling is sequential (slow for high-dimensional objects like images).
\end{frame}

\begin{frame}{Variational autoencoders (VAE)}
Latent variable model:
\[
p_\theta(x)=\int p_\theta(x\mid z)p(z)\,dz.
\]
Introduce amortized posterior \(q_\phi(z\mid x)\) and maximize ELBO:
\[
\log p_\theta(x)\ge \mathbb{E}_{q_\phi(z\mid x)}[\log p_\theta(x\mid z)]-\mathrm{KL}(q_\phi(z\mid x)\|p(z)).
\]
\end{frame}

\begin{frame}{Normalizing flows (NF)}
Invertible mapping \(x=f_\theta(z)\), \(z\sim p(z)\).
Change of variables:
\[
\log p_\theta(x)=\log p(z)+\log\left|\det \frac{\partial f_\theta^{-1}(x)}{\partial x}\right|.
\]
Pros: exact likelihood + fast sampling.
Cons: invertibility constraints and Jacobian computation structure.
\end{frame}

\begin{frame}{Energy-based models (EBM)}
Unnormalized density:
\[
p_\theta(x)=\frac{e^{-E_\theta(x)}}{Z(\theta)},\quad Z(\theta)=\int e^{-E_\theta(x)}dx.
\]
Pros: flexible.
Cons: partition function and sampling (often MCMC) are hard.
Score-based learning can avoid \(Z(\theta)\) by learning \(\nabla_x\log p\).
\end{frame}

\begin{frame}{GANs}
Implicit model: \(x=G_\theta(z)\), \(z\sim p(z)\), trained via adversarial game.
Pros: high-fidelity samples.
Cons: instability, mode collapse, no likelihood.
\end{frame}

\begin{frame}{Where diffusion fits}
Diffusion models define:
\begin{itemize}
\item a forward \emph{noising} process (easy),
\item a learned reverse \emph{denoising} process (hard, learned).
\end{itemize}
Sampling: integrate reverse dynamics from noise to data.
Training often reduces to regression on noisy inputs with time conditioning.
\end{frame}

\begin{frame}{Chapter 1 summary}
\begin{itemize}
\item Generative modeling has multiple objective families.
\item Diffusion models are dynamical/iterative generators.
\item Next: variational diffusion (DDPM) as a hierarchical latent-variable model.
\end{itemize}
\end{frame}

% =========================================================
\section{Chapter 2: Variational View (DDPM)}
% =========================================================

\begin{frame}{Diffusion as hierarchical latent model}
Treat the noised variables as latents:
\[
x_0 \sim p_{\text{data}},\quad x_1,\dots,x_T\ \text{latent},\quad x_T\approx \mathcal{N}(0,I).
\]
Forward Markov chain \(q\) is fixed; reverse chain \(p_\theta\) is learned.
\end{frame}

\begin{frame}{Forward noising process}
Common discrete forward transitions:
\[
q(x_t\mid x_{t-1})=\mathcal{N}\!\left(\sqrt{1-\beta_t}\,x_{t-1},\ \beta_t I\right),
\]
with noise schedule \(\{\beta_t\}\) and \(\alpha_t=1-\beta_t\).
\end{frame}

\begin{frame}{Closed form of the marginal \(q(x_t\mid x_0)\)}
By repeated composition,
\[
q(x_t\mid x_0)=\mathcal{N}\!\left(\sqrt{\bar\alpha_t}\,x_0,\ (1-\bar\alpha_t)I\right),
\qquad \bar\alpha_t=\prod_{s=1}^t \alpha_s.
\]
Equivalently:
\[
x_t=\sqrt{\bar\alpha_t}\,x_0+\sqrt{1-\bar\alpha_t}\,\epsilon,\quad \epsilon\sim \mathcal{N}(0,I).
\]
\end{frame}

\begin{frame}{Reverse model}
Learn a reverse Markov chain:
\[
p_\theta(x_{t-1}\mid x_t)=\mathcal{N}(\mu_\theta(x_t,t),\Sigma_\theta(x_t,t)).
\]
Often \(\Sigma_\theta\) is fixed/parameterized simply; the network predicts \(\mu_\theta\) (or an equivalent target).
\end{frame}

\begin{frame}{ELBO: variational lower bound}
Define joint:
\[
p_\theta(x_{0:T}) = p(x_T)\prod_{t=1}^T p_\theta(x_{t-1}\mid x_t),
\quad
q(x_{1:T}\mid x_0)=\prod_{t=1}^T q(x_t\mid x_{t-1}).
\]
Then
\[
\log p_\theta(x_0) \ge \mathbb{E}_{q(x_{1:T}\mid x_0)}\left[\log \frac{p_\theta(x_{0:T})}{q(x_{1:T}\mid x_0)}\right].
\]
\end{frame}

\begin{frame}{ELBO decomposition (structure)}
The bound decomposes into:
\begin{itemize}
\item a prior term \(\mathrm{KL}(q(x_T\mid x_0)\|p(x_T))\),
\item denoising KL terms \(\sum_{t=2}^T \mathrm{KL}(q(x_{t-1}\mid x_t,x_0)\|p_\theta(x_{t-1}\mid x_t))\),
\item and a reconstruction term at \(t=1\).
\end{itemize}
Because all relevant distributions are Gaussian, these KLs are tractable.
\end{frame}

\begin{frame}{Posterior \(q(x_{t-1}\mid x_t,x_0)\)}
The forward model implies an analytic posterior:
\[
q(x_{t-1}\mid x_t,x_0)=\mathcal{N}(\tilde\mu(x_t,x_0,t),\tilde\beta_t I),
\]
with closed-form \(\tilde\mu,\tilde\beta_t\) in terms of \(\alpha_t,\bar\alpha_t\).
This turns the KL term into a squared error between means (up to weights).
\end{frame}

\begin{frame}{Noise prediction parameterization}
A common choice is to predict the noise \(\epsilon\) in
\[
x_t=\sqrt{\bar\alpha_t}x_0+\sqrt{1-\bar\alpha_t}\epsilon.
\]
Let \(\epsilon_\theta(x_t,t)\) be the network. Then training becomes:
\[
\mathcal{L}_{\text{simple}}(\theta)=\mathbb{E}_{t,x_0,\epsilon}\left[\|\epsilon-\epsilon_\theta(x_t,t)\|^2\right].
\]
\end{frame}

\begin{frame}{From \(\epsilon\)-prediction to \(x_0\)-prediction}
Given \(\epsilon_\theta\), estimate the clean sample:
\[
\hat x_{0,\theta}(x_t,t)=\frac{1}{\sqrt{\bar\alpha_t}}\left(x_t-\sqrt{1-\bar\alpha_t}\,\epsilon_\theta(x_t,t)\right).
\]
This reparameterization is useful for alternative losses and samplers.
\end{frame}

\begin{frame}{Alternative targets: score and velocity}
Other equivalent prediction targets:
\begin{itemize}
\item score: \(s_\theta(x_t,t)\approx \nabla_{x_t}\log p_t(x_t)\),
\item velocity (``\(v\)''-prediction): a linear combination of \(x_t\) and \(\epsilon\),
\item direct mean prediction \(\mu_\theta\).
\end{itemize}
Chapter 6 will emphasize these are different parameterizations of a common vector field.
\end{frame}

\begin{frame}{Sampling: ancestral reverse chain}
Start with \(x_T\sim\mathcal{N}(0,I)\).
For \(t=T,\dots,1\) sample
\[
x_{t-1}\sim p_\theta(x_{t-1}\mid x_t).
\]
Each step reduces noise and adds structure; after \(T\) steps obtain \(x_0\sim p_\theta\).
\end{frame}

\begin{frame}{Discrete acceleration idea (preview)}
Discrete reverse chain is a particular discretization.
Later chapters typically discuss:
\begin{itemize}
\item fewer steps (DDIM / deterministic paths),
\item higher-order solvers,
\item distillation/consistency models.
\end{itemize}
But the foundation is the same: integrate reverse dynamics.
\end{frame}

\begin{frame}{Chapter 2 summary}
\begin{itemize}
\item Forward diffusion is fixed and Gaussian.
\item Reverse process is learned.
\item ELBO leads to tractable training; in practice, a simple MSE loss works well.
\item Next: score matching and denoising score matching.
\end{itemize}
\end{frame}

% =========================================================
\section{Chapter 3: Score-Based View (Score Matching, DSM, NCSN)}
% =========================================================

\begin{frame}{Score function}
For a density \(p(x)\), the score is
\[
s(x)\equiv \nabla_x \log p(x).
\]
Scores define directions of increasing log-density.
If we can estimate \(s\), we can build samplers that move noise toward data modes.
\end{frame}

\begin{frame}{Why score matching?}
For EBMs \(p_\theta(x)\propto e^{-E_\theta(x)}\), likelihood training requires \(Z(\theta)\).
Score matching can avoid explicit normalization by matching \(\nabla_x\log p\).
\end{frame}

\begin{frame}{Classical score matching (idea)}
Objective (conceptually):
\[
\min_\theta \ \mathbb{E}_{p_{\text{data}}}\left[\|s_\theta(x)-\nabla_x\log p_{\text{data}}(x)\|^2\right],
\]
but \(\nabla_x\log p_{\text{data}}\) is unknown.
Hyv\"arinen's reformulation yields an equivalent objective without the unknown score (under regularity).
\end{frame}

\begin{frame}{Denoising score matching (DSM)}
Corrupt data with Gaussian noise:
\[
\tilde x = x + \sigma \epsilon,\quad \epsilon\sim \mathcal{N}(0,I).
\]
Train \(s_\theta(\tilde x,\sigma)\) to approximate the score of the noisy distribution.
A regression target exists:
\[
\nabla_{\tilde x}\log p(\tilde x\mid x) = -\frac{\tilde x-x}{\sigma^2}.
\]
\end{frame}

\begin{frame}{DSM loss (regression form)}
One common DSM loss:
\[
\mathcal{L}_{\text{DSM}}(\theta)=\mathbb{E}_{x,\epsilon}\left[\left\|s_\theta(x+\sigma\epsilon,\sigma)+\frac{1}{\sigma}\epsilon\right\|^2\right],
\]
up to weighting conventions.
Key: conditioning on \(\sigma\) allows learning scores across noise scales.
\end{frame}

\begin{frame}{Noise-Conditional Score Networks (NCSN)}
Train a single network with noise conditioning:
\[
s_\theta(\tilde x,\sigma_k),\quad \sigma_k\ \text{from a noise ladder}.
\]
At large \(\sigma\), density is smooth; at small \(\sigma\), fine details matter.
This mirrors coarse-to-fine generation.
\end{frame}

\begin{frame}{Sampling: Langevin dynamics}
Given a score estimate \(s_\theta\), Langevin sampling:
\[
x^{(n+1)} = x^{(n)} + \eta\, s_\theta(x^{(n)},\sigma) + \sqrt{2\eta}\,\xi^{(n)},\quad \xi^{(n)}\sim\mathcal{N}(0,I).
\]
Annealed Langevin: progressively decrease \(\sigma\) and step sizes.
\end{frame}

\begin{frame}{Connection to diffusion denoisers}
In Gaussian noising, score and noise prediction relate roughly by:
\[
s_\theta(x_t,t) \approx -\frac{1}{\sigma_t}\,\epsilon_\theta(x_t,t)
\]
(modulo scaling conventions).
Continuous-time formulations make these equivalences exact.
\end{frame}

\begin{frame}{Chapter 3 summary}
\begin{itemize}
\item Score-based learning targets \(\nabla\log p_t\).
\item DSM provides a tractable regression target by corrupting data.
\item Sampling can be done with annealed Langevin dynamics.
\item Next: continuous-time SDE formulation.
\end{itemize}
\end{frame}

% =========================================================
\section{Chapter 4: Continuous-Time Formulation (Score SDE)}
% =========================================================

\begin{frame}{Why continuous time?}
Continuous-time view:
\begin{itemize}
\item unifies DDPM discretizations and NCSN noise ladders,
\item leverages SDE/ODE solvers,
\item gives principled reverse-time dynamics.
\end{itemize}
\end{frame}

\begin{frame}{Forward SDE}
A generic forward SDE:
\[
dx = f(x,t)\,dt + g(t)\,dw_t,
\]
where \(w_t\) is Brownian motion.
The marginal density \(p_t(x)\) evolves smoothly with time.
\end{frame}

\begin{frame}{Fokker--Planck equation}
The density evolution satisfies a PDE of the form:
\[
\partial_t p_t(x)= -\nabla\cdot(f(x,t)p_t(x))+\frac12\nabla\cdot\left(\nabla(g(t)^2 p_t(x))\right),
\]
(with isotropic diffusion).
This is the stochastic analogue of a continuity/transport equation.
\end{frame}

\begin{frame}{Reverse-time SDE (key idea)}
Reverse-time dynamics (informally) has drift involving the score:
\[
dx = \left(f(x,t)-g(t)^2 \nabla_x\log p_t(x)\right)\,dt + g(t)\,d\bar w_t,
\]
integrated from \(t=T\) back to \(t=0\).
Thus generation requires an estimator \(s_\theta(x,t)\approx \nabla\log p_t(x)\).
\end{frame}

\begin{frame}{Training in continuous time}
We learn \(s_\theta(x,t)\) by a time-conditioned denoising score matching objective.
Discrete-time DDPM becomes a particular discretization of the continuous-time objective (and vice versa).
\end{frame}

\begin{frame}{Probability flow ODE}
There exists a deterministic ODE sharing the same marginal densities:
\[
\frac{dx}{dt} = f(x,t)-\frac12 g(t)^2 s_\theta(x,t).
\]
Sampling can be done by solving this ODE backward from noise to data.
This yields deterministic samplers (related to DDIM) and enables advanced ODE solvers.
\end{frame}

\begin{frame}{Samplers as numerical integrators}
Sampling reduces to integrating an SDE or ODE:
\begin{itemize}
\item DDPM ancestral sampling: Euler--Maruyama-style discretization of reverse SDE.
\item Deterministic samplers: discretizations of PF-ODE.
\item Higher-order solvers can trade compute vs fidelity.
\end{itemize}
\end{frame}

\begin{frame}{Chapter 4 summary}
\begin{itemize}
\item SDE formulation unifies diffusion training/sampling.
\item Reverse-time requires the score.
\item PF-ODE provides deterministic sampling with identical marginals.
\item Next: flow-based view and flow matching.
\end{itemize}
\end{frame}

% =========================================================
\section{Chapter 5: Flow-Based View (NF, CNF/Neural ODE, Flow Matching)}
% =========================================================

\begin{frame}{Normalizing flows revisited}
Flows transport a simple prior to the data:
\[
z\sim p_0 \xrightarrow{x=f_\theta(z)} x\sim p_\theta.
\]
If \(f_\theta\) is invertible, likelihood is exact via change of variables.
\end{frame}

\begin{frame}{Continuous normalizing flows (CNF)}
Define an ODE:
\[
\frac{dz}{dt}=v_\phi(z(t),t),\quad z(0)\sim p_0,\quad x=z(1).
\]
Density evolves by:
\[
\frac{d}{dt}\log p(z(t),t) = -\nabla\cdot v_\phi(z(t),t).
\]
Thus likelihood requires integrating divergence along paths.
\end{frame}

\begin{frame}{Training and the adjoint}
Differentiating through ODE solutions can use the adjoint method:
\begin{itemize}
\item low memory footprint,
\item but requires solving auxiliary ODEs backward,
\item can be computationally expensive for high-dimensional data.
\end{itemize}
This motivates alternative training objectives that avoid likelihood integration.
\end{frame}

\begin{frame}{Flow matching: core idea}
Flow matching directly trains a vector field \(v_\phi(x,t)\) by regression:
\[
\min_\phi \ \mathbb{E}_{t,x_t}\big[\|v_\phi(x_t,t)-u_t(x_t,t)\|^2\big],
\]
where \(u_t\) is a \emph{target velocity field} defined from a chosen probability path \(p_t\).
Key: ``simulation-free'' training in many constructions.
\end{frame}

\begin{frame}{Why flow matching matters for diffusion}
\begin{itemize}
\item PF-ODE sampling already interprets diffusion as an ODE flow.
\item Flow matching formalizes learning the ODE drift directly.
\item Provides bridges to optimal transport, Schr\"odinger bridges, and continuity equations.
\end{itemize}
\end{frame}

\begin{frame}{Chapter 5 summary}
\begin{itemize}
\item Flows provide exact likelihood under invertibility.
\item CNFs extend flows to continuous time with divergence dynamics.
\item Flow matching trains vector fields by regression, aligning with diffusion's ODE view.
\item Next: unify perspectives via conditioning and PDEs.
\end{itemize}
\end{frame}

% =========================================================
\section{Chapter 6: Unifying Principles}
% =========================================================

\begin{frame}{Unifying message}
Across DDPM, score-based models, and flow matching:
\[
\textbf{learn a time-dependent vector field that transports }p_0 \to p_1.
\]
Different parameterizations (noise, \(x_0\), score, velocity) are often algebraically convertible.
\end{frame}

\begin{frame}{Conditional strategy}
Training becomes tractable by conditioning on a partially transformed sample:
\[
\text{learn}\quad \mathbb{E}[\cdot\mid x_t,t]
\]
using a neural network that takes \((x_t,t)\) (and optionally condition \(c\)) as input.
This is the backbone of:
\begin{itemize}
\item DDPM regression,
\item DSM/NCSN score learning,
\item flow matching regression.
\end{itemize}
\end{frame}

\begin{frame}{Deterministic vs stochastic descriptions}
Two descriptions that can share the same marginals:
\begin{itemize}
\item stochastic: forward SDE + reverse SDE (score in drift),
\item deterministic: probability flow ODE (vector field in ODE).
\end{itemize}
Thus ``diffusion sampling'' is, fundamentally, \emph{solving a differential equation}.
\end{frame}

\begin{frame}{Continuity equation for deterministic flows}
For ODE \( \dot x = v(x,t)\), density satisfies:
\[
\partial_t p_t(x) = -\nabla\cdot(p_t(x)\,v(x,t)).
\]
This is a conservation law for probability mass.
\end{frame}

\begin{frame}{Fokker--Planck for SDEs}
For SDE \(dx=f\,dt+g\,dw\), density evolves by:
\[
\partial_t p_t = -\nabla\cdot(f p_t) + \frac12 \nabla\cdot(\nabla(g^2 p_t)).
\]
Diffusion adds a smoothing (second-order) term.
The reverse-time drift involves the score to counteract this smoothing during generation.
\end{frame}

\begin{frame}{Parameterizations as re-labelings}
Common choices:
\begin{itemize}
\item \(\epsilon\)-prediction: predict added noise.
\item \(x_0\)-prediction: predict clean data (denoising).
\item score: \(s_\theta(x,t)\approx \nabla\log p_t(x)\).
\item velocity: \(v_\theta(x,t)\) in PF-ODE/flow matching.
\end{itemize}
Conversions depend on the forward process (VP/VE/sub-VP) and scaling.
\end{frame}

\begin{frame}{Practical implications}
\begin{itemize}
\item Training stability depends on parameterization and weighting (often SNR-based).
\item Sampling quality and speed depend on discretization/solver.
\item Unifying view suggests ``algorithm design'' is numerical analysis of differential equations.
\end{itemize}
\end{frame}

\begin{frame}{Chapter 6 summary}
\begin{itemize}
\item Conditioning + differential equations unify diffusion methodologies.
\item DDPM/score/flow matching differ mainly in how they train and discretize the same underlying object.
\end{itemize}
\end{frame}

% =========================================================
\section{Recitations, Mini-Exercises, and Slide Fillers (to reach 100+)}
% =========================================================

\begin{frame}{Mini-exercises: Chapter 1}
\begin{enumerate}
\item Compare AR, VAE, NF, EBM, GAN: objective, sampling speed, mode coverage.
\item Why does \(\mathrm{KL}(p_{\text{data}}\|p_\theta)\) penalize missing modes strongly?
\item Give an example where likelihood improves but perceptual quality worsens.
\end{enumerate}
\end{frame}

\begin{frame}{Mini-exercises: Chapter 2}
\begin{enumerate}
\item Derive \(q(x_t\mid x_0)\) by repeated substitution.
\item From \(\epsilon\)-prediction, derive \(\hat x_0\) and \(\hat\mu_\theta\) used in the reverse Gaussian.
\item Discuss why predicting \(\epsilon\) is numerically stable at large \(t\).
\end{enumerate}
\end{frame}

\begin{frame}{Mini-exercises: Chapter 3}
\begin{enumerate}
\item Derive DSM target for Gaussian corruption.
\item Explain annealed Langevin: why do we reduce \(\sigma\) gradually?
\item Compare Langevin sampling vs reverse diffusion sampling conceptually.
\end{enumerate}
\end{frame}

\begin{frame}{Mini-exercises: Chapter 4}
\begin{enumerate}
\item State the general reverse-time SDE drift and interpret the score term.
\item Show that PF-ODE and reverse SDE share the same marginals (outline).
\item Discuss numerical stability of Euler--Maruyama vs higher-order solvers.
\end{enumerate}
\end{frame}

\begin{frame}{Mini-exercises: Chapter 5}
\begin{enumerate}
\item For ODE \( \dot z=v(z,t)\), derive \( \frac{d}{dt}\log p(z(t),t)=-\nabla\cdot v\).
\item Why is divergence computation hard in high dimensions? Name a stochastic estimator.
\item Construct a simple linear flow matching example (Gaussian to Gaussian).
\end{enumerate}
\end{frame}

\begin{frame}{Mini-exercises: Chapter 6}
\begin{enumerate}
\item Write continuity and Fokker--Planck equations and explain the difference.
\item Give a mapping between \(\epsilon\)-prediction and score for a Gaussian noising process.
\item Explain why ``learning a vector field'' is the unifying principle.
\end{enumerate}
\end{frame}

% --- Add many short recap frames to exceed 100 frames ---
\begin{frame}{Recap 1: Key equations}
\begin{itemize}
\item Forward diffusion: \(q(x_t\mid x_{t-1})\).
\item Closed-form noising: \(x_t=\sqrt{\bar\alpha_t}x_0+\sqrt{1-\bar\alpha_t}\epsilon\).
\item Reverse-time drift uses score: \(f-g^2 \nabla\log p_t\).
\item PF-ODE: \(\dot x=f-\tfrac12 g^2 s_\theta\).
\end{itemize}
\end{frame}

\begin{frame}{Recap 2: What is learned?}
\begin{itemize}
\item DDPM: \(\epsilon_\theta(x_t,t)\) (or \(\hat x_0\), \(\mu_\theta\)).
\item Score methods: \(s_\theta(x,t)\approx \nabla\log p_t(x)\).
\item Flow matching: \(v_\phi(x,t)\) (velocity field).
\end{itemize}
\end{frame}

\begin{frame}{Recap 3: Sampling viewpoint}
\begin{itemize}
\item Discrete-time: iterate \(x_T\to x_0\) by learned transitions.
\item Continuous-time: integrate reverse SDE or PF-ODE.
\item All are numerical schemes for differential equations.
\end{itemize}
\end{frame}

\begin{frame}{Recap 4: Conditioning}
\begin{itemize}
\item Time conditioning is essential: model must know noise level.
\item Conditional generation adds \(c\) (labels, text, measurements).
\item Guidance modifies the vector field/score to steer samples.
\end{itemize}
\end{frame}

% Generate a batch of short "concept check" frames to bring total > 100

\begin{frame}{Concept check: KL directions}
\begin{block}{Question}
Why does \(\mathrm{KL}(p_{\text{data}}\|p_\theta)\) penalize missing modes more than \(\mathrm{KL}(p_\theta\|p_{\text{data}})\)?
\end{block}
\begin{block}{Discussion prompts}
\begin{itemize}
\item Relate to the chapter perspective (variational / score / flow).
\item Write down the relevant equation and identify learned vs fixed terms.
\item Consider numerical implications for training or sampling.
\end{itemize}
\end{block}
\end{frame}

\begin{frame}{Concept check: Markov property}
\begin{block}{Question}
In DDPM forward process, why is \(q(x_t\mid x_{t-1})\) chosen Markovian?
\end{block}
\begin{block}{Discussion prompts}
\begin{itemize}
\item Relate to the chapter perspective (variational / score / flow).
\item Write down the relevant equation and identify learned vs fixed terms.
\item Consider numerical implications for training or sampling.
\end{itemize}
\end{block}
\end{frame}

\begin{frame}{Concept check: Gaussianity}
\begin{block}{Question}
Which assumptions make \(q(x_t\mid x_0)\) Gaussian?
\end{block}
\begin{block}{Discussion prompts}
\begin{itemize}
\item Relate to the chapter perspective (variational / score / flow).
\item Write down the relevant equation and identify learned vs fixed terms.
\item Consider numerical implications for training or sampling.
\end{itemize}
\end{block}
\end{frame}

\begin{frame}{Concept check: regression view}
\begin{block}{Question}
Why does the ELBO reduce to an MSE loss under \(\epsilon\)-prediction?
\end{block}
\begin{block}{Discussion prompts}
\begin{itemize}
\item Relate to the chapter perspective (variational / score / flow).
\item Write down the relevant equation and identify learned vs fixed terms.
\item Consider numerical implications for training or sampling.
\end{itemize}
\end{block}
\end{frame}

\begin{frame}{Concept check: score meaning}
\begin{block}{Question}
Interpret \(s(x)=\nabla_x\log p(x)\) geometrically.
\end{block}
\begin{block}{Discussion prompts}
\begin{itemize}
\item Relate to the chapter perspective (variational / score / flow).
\item Write down the relevant equation and identify learned vs fixed terms.
\item Consider numerical implications for training or sampling.
\end{itemize}
\end{block}
\end{frame}

\begin{frame}{Concept check: Langevin}
\begin{block}{Question}
Why does Langevin dynamics sample from \(p(x)\) when using the exact score?
\end{block}
\begin{block}{Discussion prompts}
\begin{itemize}
\item Relate to the chapter perspective (variational / score / flow).
\item Write down the relevant equation and identify learned vs fixed terms.
\item Consider numerical implications for training or sampling.
\end{itemize}
\end{block}
\end{frame}

\begin{frame}{Concept check: annealing}
\begin{block}{Question}
What is the role of decreasing \(\sigma\) in annealed Langevin dynamics?
\end{block}
\begin{block}{Discussion prompts}
\begin{itemize}
\item Relate to the chapter perspective (variational / score / flow).
\item Write down the relevant equation and identify learned vs fixed terms.
\item Consider numerical implications for training or sampling.
\end{itemize}
\end{block}
\end{frame}

\begin{frame}{Concept check: Fokker--Planck}
\begin{block}{Question}
What does the second-order term in Fokker--Planck represent physically?
\end{block}
\begin{block}{Discussion prompts}
\begin{itemize}
\item Relate to the chapter perspective (variational / score / flow).
\item Write down the relevant equation and identify learned vs fixed terms.
\item Consider numerical implications for training or sampling.
\end{itemize}
\end{block}
\end{frame}

\begin{frame}{Concept check: reverse drift}
\begin{block}{Question}
Why must the reverse drift include the score term?
\end{block}
\begin{block}{Discussion prompts}
\begin{itemize}
\item Relate to the chapter perspective (variational / score / flow).
\item Write down the relevant equation and identify learned vs fixed terms.
\item Consider numerical implications for training or sampling.
\end{itemize}
\end{block}
\end{frame}

\begin{frame}{Concept check: PF-ODE}
\begin{block}{Question}
Why can an ODE share marginals with an SDE?
\end{block}
\begin{block}{Discussion prompts}
\begin{itemize}
\item Relate to the chapter perspective (variational / score / flow).
\item Write down the relevant equation and identify learned vs fixed terms.
\item Consider numerical implications for training or sampling.
\end{itemize}
\end{block}
\end{frame}

\begin{frame}{Concept check: divergence}
\begin{block}{Question}
In CNFs, why does divergence appear in the log-density dynamics?
\end{block}
\begin{block}{Discussion prompts}
\begin{itemize}
\item Relate to the chapter perspective (variational / score / flow).
\item Write down the relevant equation and identify learned vs fixed terms.
\item Consider numerical implications for training or sampling.
\end{itemize}
\end{block}
\end{frame}

\begin{frame}{Concept check: flow matching}
\begin{block}{Question}
What does flow matching regress and why can it avoid likelihood evaluation?
\end{block}
\begin{block}{Discussion prompts}
\begin{itemize}
\item Relate to the chapter perspective (variational / score / flow).
\item Write down the relevant equation and identify learned vs fixed terms.
\item Consider numerical implications for training or sampling.
\end{itemize}
\end{block}
\end{frame}

\begin{frame}{Concept check: parameterizations}
\begin{block}{Question}
Give one practical advantage of \(\epsilon\)-prediction vs \(x_0\)-prediction.
\end{block}
\begin{block}{Discussion prompts}
\begin{itemize}
\item Relate to the chapter perspective (variational / score / flow).
\item Write down the relevant equation and identify learned vs fixed terms.
\item Consider numerical implications for training or sampling.
\end{itemize}
\end{block}
\end{frame}

\begin{frame}{Concept check: discretization}
\begin{block}{Question}
How does step size affect sampling quality and computational cost?
\end{block}
\begin{block}{Discussion prompts}
\begin{itemize}
\item Relate to the chapter perspective (variational / score / flow).
\item Write down the relevant equation and identify learned vs fixed terms.
\item Consider numerical implications for training or sampling.
\end{itemize}
\end{block}
\end{frame}

\begin{frame}{Concept check: conditioning}
\begin{block}{Question}
Why is time conditioning essential for a single network across all noise levels?
\end{block}
\begin{block}{Discussion prompts}
\begin{itemize}
\item Relate to the chapter perspective (variational / score / flow).
\item Write down the relevant equation and identify learned vs fixed terms.
\item Consider numerical implications for training or sampling.
\end{itemize}
\end{block}
\end{frame}

\begin{frame}{Concept check 16: KL directions}
\begin{block}{Question}
Why does \(\mathrm{KL}(p_{\text{data}}\|p_\theta)\) penalize missing modes more than \(\mathrm{KL}(p_\theta\|p_{\text{data}})\)?
\end{block}
\begin{block}{Discussion prompts}
\begin{itemize}
\item Relate to the chapter perspective (variational / score / flow).
\item Write down the relevant equation and identify learned vs fixed terms.
\item Consider numerical implications for training or sampling.
\end{itemize}
\end{block}
\end{frame}

\begin{frame}{Concept check 17: Markov property}
\begin{block}{Question}
In DDPM forward process, why is \(q(x_t\mid x_{t-1})\) chosen Markovian?
\end{block}
\begin{block}{Discussion prompts}
\begin{itemize}
\item Relate to the chapter perspective (variational / score / flow).
\item Write down the relevant equation and identify learned vs fixed terms.
\item Consider numerical implications for training or sampling.
\end{itemize}
\end{block}
\end{frame}

\begin{frame}{Concept check 18: Gaussianity}
\begin{block}{Question}
Which assumptions make \(q(x_t\mid x_0)\) Gaussian?
\end{block}
\begin{block}{Discussion prompts}
\begin{itemize}
\item Relate to the chapter perspective (variational / score / flow).
\item Write down the relevant equation and identify learned vs fixed terms.
\item Consider numerical implications for training or sampling.
\end{itemize}
\end{block}
\end{frame}

\begin{frame}{Concept check 19: regression view}
\begin{block}{Question}
Why does the ELBO reduce to an MSE loss under \(\epsilon\)-prediction?
\end{block}
\begin{block}{Discussion prompts}
\begin{itemize}
\item Relate to the chapter perspective (variational / score / flow).
\item Write down the relevant equation and identify learned vs fixed terms.
\item Consider numerical implications for training or sampling.
\end{itemize}
\end{block}
\end{frame}

\begin{frame}{Concept check 20: score meaning}
\begin{block}{Question}
Interpret \(s(x)=\nabla_x\log p(x)\) geometrically.
\end{block}
\begin{block}{Discussion prompts}
\begin{itemize}
\item Relate to the chapter perspective (variational / score / flow).
\item Write down the relevant equation and identify learned vs fixed terms.
\item Consider numerical implications for training or sampling.
\end{itemize}
\end{block}
\end{frame}

\begin{frame}{Concept check 21: Langevin}
\begin{block}{Question}
Why does Langevin dynamics sample from \(p(x)\) when using the exact score?
\end{block}
\begin{block}{Discussion prompts}
\begin{itemize}
\item Relate to the chapter perspective (variational / score / flow).
\item Write down the relevant equation and identify learned vs fixed terms.
\item Consider numerical implications for training or sampling.
\end{itemize}
\end{block}
\end{frame}

\begin{frame}{Concept check 22: annealing}
\begin{block}{Question}
What is the role of decreasing \(\sigma\) in annealed Langevin dynamics?
\end{block}
\begin{block}{Discussion prompts}
\begin{itemize}
\item Relate to the chapter perspective (variational / score / flow).
\item Write down the relevant equation and identify learned vs fixed terms.
\item Consider numerical implications for training or sampling.
\end{itemize}
\end{block}
\end{frame}

\begin{frame}{Concept check 23: Fokker--Planck}
\begin{block}{Question}
What does the second-order term in Fokker--Planck represent physically?
\end{block}
\begin{block}{Discussion prompts}
\begin{itemize}
\item Relate to the chapter perspective (variational / score / flow).
\item Write down the relevant equation and identify learned vs fixed terms.
\item Consider numerical implications for training or sampling.
\end{itemize}
\end{block}
\end{frame}

\begin{frame}{Concept check 24: reverse drift}
\begin{block}{Question}
Why must the reverse drift include the score term?
\end{block}
\begin{block}{Discussion prompts}
\begin{itemize}
\item Relate to the chapter perspective (variational / score / flow).
\item Write down the relevant equation and identify learned vs fixed terms.
\item Consider numerical implications for training or sampling.
\end{itemize}
\end{block}
\end{frame}

\begin{frame}{Concept check 25: PF-ODE}
\begin{block}{Question}
Why can an ODE share marginals with an SDE?
\end{block}
\begin{block}{Discussion prompts}
\begin{itemize}
\item Relate to the chapter perspective (variational / score / flow).
\item Write down the relevant equation and identify learned vs fixed terms.
\item Consider numerical implications for training or sampling.
\end{itemize}
\end{block}
\end{frame}

\begin{frame}{Concept check 26: divergence}
\begin{block}{Question}
In CNFs, why does divergence appear in the log-density dynamics?
\end{block}
\begin{block}{Discussion prompts}
\begin{itemize}
\item Relate to the chapter perspective (variational / score / flow).
\item Write down the relevant equation and identify learned vs fixed terms.
\item Consider numerical implications for training or sampling.
\end{itemize}
\end{block}
\end{frame}

\begin{frame}{Concept check 27: flow matching}
\begin{block}{Question}
What does flow matching regress and why can it avoid likelihood evaluation?
\end{block}
\begin{block}{Discussion prompts}
\begin{itemize}
\item Relate to the chapter perspective (variational / score / flow).
\item Write down the relevant equation and identify learned vs fixed terms.
\item Consider numerical implications for training or sampling.
\end{itemize}
\end{block}
\end{frame}

\begin{frame}{Concept check 28: parameterizations}
\begin{block}{Question}
Give one practical advantage of \(\epsilon\)-prediction vs \(x_0\)-prediction.
\end{block}
\begin{block}{Discussion prompts}
\begin{itemize}
\item Relate to the chapter perspective (variational / score / flow).
\item Write down the relevant equation and identify learned vs fixed terms.
\item Consider numerical implications for training or sampling.
\end{itemize}
\end{block}
\end{frame}

\begin{frame}{Concept check 29: discretization}
\begin{block}{Question}
How does step size affect sampling quality and computational cost?
\end{block}
\begin{block}{Discussion prompts}
\begin{itemize}
\item Relate to the chapter perspective (variational / score / flow).
\item Write down the relevant equation and identify learned vs fixed terms.
\item Consider numerical implications for training or sampling.
\end{itemize}
\end{block}
\end{frame}

\begin{frame}{Concept check 30: conditioning}
\begin{block}{Question}
Why is time conditioning essential for a single network across all noise levels?
\end{block}
\begin{block}{Discussion prompts}
\begin{itemize}
\item Relate to the chapter perspective (variational / score / flow).
\item Write down the relevant equation and identify learned vs fixed terms.
\item Consider numerical implications for training or sampling.
\end{itemize}
\end{block}
\end{frame}

\begin{frame}{Concept check 31: KL directions}
\begin{block}{Question}
Why does \(\mathrm{KL}(p_{\text{data}}\|p_\theta)\) penalize missing modes more than \(\mathrm{KL}(p_\theta\|p_{\text{data}})\)?
\end{block}
\begin{block}{Discussion prompts}
\begin{itemize}
\item Relate to the chapter perspective (variational / score / flow).
\item Write down the relevant equation and identify learned vs fixed terms.
\item Consider numerical implications for training or sampling.
\end{itemize}
\end{block}
\end{frame}

\begin{frame}{Concept check 32: Markov property}
\begin{block}{Question}
In DDPM forward process, why is \(q(x_t\mid x_{t-1})\) chosen Markovian?
\end{block}
\begin{block}{Discussion prompts}
\begin{itemize}
\item Relate to the chapter perspective (variational / score / flow).
\item Write down the relevant equation and identify learned vs fixed terms.
\item Consider numerical implications for training or sampling.
\end{itemize}
\end{block}
\end{frame}

\begin{frame}{Concept check 33: Gaussianity}
\begin{block}{Question}
Which assumptions make \(q(x_t\mid x_0)\) Gaussian?
\end{block}
\begin{block}{Discussion prompts}
\begin{itemize}
\item Relate to the chapter perspective (variational / score / flow).
\item Write down the relevant equation and identify learned vs fixed terms.
\item Consider numerical implications for training or sampling.
\end{itemize}
\end{block}
\end{frame}

\begin{frame}{Concept check 34: regression view}
\begin{block}{Question}
Why does the ELBO reduce to an MSE loss under \(\epsilon\)-prediction?
\end{block}
\begin{block}{Discussion prompts}
\begin{itemize}
\item Relate to the chapter perspective (variational / score / flow).
\item Write down the relevant equation and identify learned vs fixed terms.
\item Consider numerical implications for training or sampling.
\end{itemize}
\end{block}
\end{frame}

\begin{frame}{Concept check 35: score meaning}
\begin{block}{Question}
Interpret \(s(x)=\nabla_x\log p(x)\) geometrically.
\end{block}
\begin{block}{Discussion prompts}
\begin{itemize}
\item Relate to the chapter perspective (variational / score / flow).
\item Write down the relevant equation and identify learned vs fixed terms.
\item Consider numerical implications for training or sampling.
\end{itemize}
\end{block}
\end{frame}

\begin{frame}{Concept check 36: Langevin}
\begin{block}{Question}
Why does Langevin dynamics sample from \(p(x)\) when using the exact score?
\end{block}
\begin{block}{Discussion prompts}
\begin{itemize}
\item Relate to the chapter perspective (variational / score / flow).
\item Write down the relevant equation and identify learned vs fixed terms.
\item Consider numerical implications for training or sampling.
\end{itemize}
\end{block}
\end{frame}

\begin{frame}{Concept check 37: annealing}
\begin{block}{Question}
What is the role of decreasing \(\sigma\) in annealed Langevin dynamics?
\end{block}
\begin{block}{Discussion prompts}
\begin{itemize}
\item Relate to the chapter perspective (variational / score / flow).
\item Write down the relevant equation and identify learned vs fixed terms.
\item Consider numerical implications for training or sampling.
\end{itemize}
\end{block}
\end{frame}

\begin{frame}{Concept check 38: Fokker--Planck}
\begin{block}{Question}
What does the second-order term in Fokker--Planck represent physically?
\end{block}
\begin{block}{Discussion prompts}
\begin{itemize}
\item Relate to the chapter perspective (variational / score / flow).
\item Write down the relevant equation and identify learned vs fixed terms.
\item Consider numerical implications for training or sampling.
\end{itemize}
\end{block}
\end{frame}

\begin{frame}{Concept check 39: reverse drift}
\begin{block}{Question}
Why must the reverse drift include the score term?
\end{block}
\begin{block}{Discussion prompts}
\begin{itemize}
\item Relate to the chapter perspective (variational / score / flow).
\item Write down the relevant equation and identify learned vs fixed terms.
\item Consider numerical implications for training or sampling.
\end{itemize}
\end{block}
\end{frame}

\begin{frame}{Concept check 40: PF-ODE}
\begin{block}{Question}
Why can an ODE share marginals with an SDE?
\end{block}
\begin{block}{Discussion prompts}
\begin{itemize}
\item Relate to the chapter perspective (variational / score / flow).
\item Write down the relevant equation and identify learned vs fixed terms.
\item Consider numerical implications for training or sampling.
\end{itemize}
\end{block}
\end{frame}

\begin{frame}{Concept check 41: divergence}
\begin{block}{Question}
In CNFs, why does divergence appear in the log-density dynamics?
\end{block}
\begin{block}{Discussion prompts}
\begin{itemize}
\item Relate to the chapter perspective (variational / score / flow).
\item Write down the relevant equation and identify learned vs fixed terms.
\item Consider numerical implications for training or sampling.
\end{itemize}
\end{block}
\end{frame}

\begin{frame}{Concept check 42: flow matching}
\begin{block}{Question}
What does flow matching regress and why can it avoid likelihood evaluation?
\end{block}
\begin{block}{Discussion prompts}
\begin{itemize}
\item Relate to the chapter perspective (variational / score / flow).
\item Write down the relevant equation and identify learned vs fixed terms.
\item Consider numerical implications for training or sampling.
\end{itemize}
\end{block}
\end{frame}

\begin{frame}{Concept check 43: parameterizations}
\begin{block}{Question}
Give one practical advantage of \(\epsilon\)-prediction vs \(x_0\)-prediction.
\end{block}
\begin{block}{Discussion prompts}
\begin{itemize}
\item Relate to the chapter perspective (variational / score / flow).
\item Write down the relevant equation and identify learned vs fixed terms.
\item Consider numerical implications for training or sampling.
\end{itemize}
\end{block}
\end{frame}

\begin{frame}{Concept check 44: discretization}
\begin{block}{Question}
How does step size affect sampling quality and computational cost?
\end{block}
\begin{block}{Discussion prompts}
\begin{itemize}
\item Relate to the chapter perspective (variational / score / flow).
\item Write down the relevant equation and identify learned vs fixed terms.
\item Consider numerical implications for training or sampling.
\end{itemize}
\end{block}
\end{frame}

\begin{frame}{Concept check 45: conditioning}
\begin{block}{Question}
Why is time conditioning essential for a single network across all noise levels?
\end{block}
\begin{block}{Discussion prompts}
\begin{itemize}
\item Relate to the chapter perspective (variational / score / flow).
\item Write down the relevant equation and identify learned vs fixed terms.
\item Consider numerical implications for training or sampling.
\end{itemize}
\end{block}
\end{frame}

\begin{frame}{Concept check 46: KL directions}
\begin{block}{Question}
Why does \(\mathrm{KL}(p_{\text{data}}\|p_\theta)\) penalize missing modes more than \(\mathrm{KL}(p_\theta\|p_{\text{data}})\)?
\end{block}
\begin{block}{Discussion prompts}
\begin{itemize}
\item Relate to the chapter perspective (variational / score / flow).
\item Write down the relevant equation and identify learned vs fixed terms.
\item Consider numerical implications for training or sampling.
\end{itemize}
\end{block}
\end{frame}

\begin{frame}{Concept check 47: Markov property}
\begin{block}{Question}
In DDPM forward process, why is \(q(x_t\mid x_{t-1})\) chosen Markovian?
\end{block}
\begin{block}{Discussion prompts}
\begin{itemize}
\item Relate to the chapter perspective (variational / score / flow).
\item Write down the relevant equation and identify learned vs fixed terms.
\item Consider numerical implications for training or sampling.
\end{itemize}
\end{block}
\end{frame}

\begin{frame}{Concept check 48: Gaussianity}
\begin{block}{Question}
Which assumptions make \(q(x_t\mid x_0)\) Gaussian?
\end{block}
\begin{block}{Discussion prompts}
\begin{itemize}
\item Relate to the chapter perspective (variational / score / flow).
\item Write down the relevant equation and identify learned vs fixed terms.
\item Consider numerical implications for training or sampling.
\end{itemize}
\end{block}
\end{frame}

\begin{frame}{Concept check 49: regression view}
\begin{block}{Question}
Why does the ELBO reduce to an MSE loss under \(\epsilon\)-prediction?
\end{block}
\begin{block}{Discussion prompts}
\begin{itemize}
\item Relate to the chapter perspective (variational / score / flow).
\item Write down the relevant equation and identify learned vs fixed terms.
\item Consider numerical implications for training or sampling.
\end{itemize}
\end{block}
\end{frame}

\begin{frame}{Concept check 50: score meaning}
\begin{block}{Question}
Interpret \(s(x)=\nabla_x\log p(x)\) geometrically.
\end{block}
\begin{block}{Discussion prompts}
\begin{itemize}
\item Relate to the chapter perspective (variational / score / flow).
\item Write down the relevant equation and identify learned vs fixed terms.
\item Consider numerical implications for training or sampling.
\end{itemize}
\end{block}
\end{frame}

\begin{frame}{Concept check 51: Langevin}
\begin{block}{Question}
Why does Langevin dynamics sample from \(p(x)\) when using the exact score?
\end{block}
\begin{block}{Discussion prompts}
\begin{itemize}
\item Relate to the chapter perspective (variational / score / flow).
\item Write down the relevant equation and identify learned vs fixed terms.
\item Consider numerical implications for training or sampling.
\end{itemize}
\end{block}
\end{frame}

\begin{frame}{Concept check 52: annealing}
\begin{block}{Question}
What is the role of decreasing \(\sigma\) in annealed Langevin dynamics?
\end{block}
\begin{block}{Discussion prompts}
\begin{itemize}
\item Relate to the chapter perspective (variational / score / flow).
\item Write down the relevant equation and identify learned vs fixed terms.
\item Consider numerical implications for training or sampling.
\end{itemize}
\end{block}
\end{frame}

\begin{frame}{Concept check 53: Fokker--Planck}
\begin{block}{Question}
What does the second-order term in Fokker--Planck represent physically?
\end{block}
\begin{block}{Discussion prompts}
\begin{itemize}
\item Relate to the chapter perspective (variational / score / flow).
\item Write down the relevant equation and identify learned vs fixed terms.
\item Consider numerical implications for training or sampling.
\end{itemize}
\end{block}
\end{frame}

\begin{frame}{Concept check 54: reverse drift}
\begin{block}{Question}
Why must the reverse drift include the score term?
\end{block}
\begin{block}{Discussion prompts}
\begin{itemize}
\item Relate to the chapter perspective (variational / score / flow).
\item Write down the relevant equation and identify learned vs fixed terms.
\item Consider numerical implications for training or sampling.
\end{itemize}
\end{block}
\end{frame}

\begin{frame}{Concept check 55: PF-ODE}
\begin{block}{Question}
Why can an ODE share marginals with an SDE?
\end{block}
\begin{block}{Discussion prompts}
\begin{itemize}
\item Relate to the chapter perspective (variational / score / flow).
\item Write down the relevant equation and identify learned vs fixed terms.
\item Consider numerical implications for training or sampling.
\end{itemize}
\end{block}
\end{frame}

\begin{frame}{Concept check 56: divergence}
\begin{block}{Question}
In CNFs, why does divergence appear in the log-density dynamics?
\end{block}
\begin{block}{Discussion prompts}
\begin{itemize}
\item Relate to the chapter perspective (variational / score / flow).
\item Write down the relevant equation and identify learned vs fixed terms.
\item Consider numerical implications for training or sampling.
\end{itemize}
\end{block}
\end{frame}

\begin{frame}{Concept check 57: flow matching}
\begin{block}{Question}
What does flow matching regress and why can it avoid likelihood evaluation?
\end{block}
\begin{block}{Discussion prompts}
\begin{itemize}
\item Relate to the chapter perspective (variational / score / flow).
\item Write down the relevant equation and identify learned vs fixed terms.
\item Consider numerical implications for training or sampling.
\end{itemize}
\end{block}
\end{frame}

\begin{frame}{Concept check 58: parameterizations}
\begin{block}{Question}
Give one practical advantage of \(\epsilon\)-prediction vs \(x_0\)-prediction.
\end{block}
\begin{block}{Discussion prompts}
\begin{itemize}
\item Relate to the chapter perspective (variational / score / flow).
\item Write down the relevant equation and identify learned vs fixed terms.
\item Consider numerical implications for training or sampling.
\end{itemize}
\end{block}
\end{frame}

\begin{frame}{Concept check 59: discretization}
\begin{block}{Question}
How does step size affect sampling quality and computational cost?
\end{block}
\begin{block}{Discussion prompts}
\begin{itemize}
\item Relate to the chapter perspective (variational / score / flow).
\item Write down the relevant equation and identify learned vs fixed terms.
\item Consider numerical implications for training or sampling.
\end{itemize}
\end{block}
\end{frame}

\begin{frame}{Concept check 60: conditioning}
\begin{block}{Question}
Why is time conditioning essential for a single network across all noise levels?
\end{block}
\begin{block}{Discussion prompts}
\begin{itemize}
\item Relate to the chapter perspective (variational / score / flow).
\item Write down the relevant equation and identify learned vs fixed terms.
\item Consider numerical implications for training or sampling.
\end{itemize}
\end{block}
\end{frame}

\begin{frame}{Concept check 61: KL directions}
\begin{block}{Question}
Why does \(\mathrm{KL}(p_{\text{data}}\|p_\theta)\) penalize missing modes more than \(\mathrm{KL}(p_\theta\|p_{\text{data}})\)?
\end{block}
\begin{block}{Discussion prompts}
\begin{itemize}
\item Relate to the chapter perspective (variational / score / flow).
\item Write down the relevant equation and identify learned vs fixed terms.
\item Consider numerical implications for training or sampling.
\end{itemize}
\end{block}
\end{frame}

\begin{frame}{Concept check 62: Markov property}
\begin{block}{Question}
In DDPM forward process, why is \(q(x_t\mid x_{t-1})\) chosen Markovian?
\end{block}
\begin{block}{Discussion prompts}
\begin{itemize}
\item Relate to the chapter perspective (variational / score / flow).
\item Write down the relevant equation and identify learned vs fixed terms.
\item Consider numerical implications for training or sampling.
\end{itemize}
\end{block}
\end{frame}

\begin{frame}{Concept check 63: Gaussianity}
\begin{block}{Question}
Which assumptions make \(q(x_t\mid x_0)\) Gaussian?
\end{block}
\begin{block}{Discussion prompts}
\begin{itemize}
\item Relate to the chapter perspective (variational / score / flow).
\item Write down the relevant equation and identify learned vs fixed terms.
\item Consider numerical implications for training or sampling.
\end{itemize}
\end{block}
\end{frame}

\begin{frame}{Concept check 64: regression view}
\begin{block}{Question}
Why does the ELBO reduce to an MSE loss under \(\epsilon\)-prediction?
\end{block}
\begin{block}{Discussion prompts}
\begin{itemize}
\item Relate to the chapter perspective (variational / score / flow).
\item Write down the relevant equation and identify learned vs fixed terms.
\item Consider numerical implications for training or sampling.
\end{itemize}
\end{block}
\end{frame}

\begin{frame}{Concept check 65: score meaning}
\begin{block}{Question}
Interpret \(s(x)=\nabla_x\log p(x)\) geometrically.
\end{block}
\begin{block}{Discussion prompts}
\begin{itemize}
\item Relate to the chapter perspective (variational / score / flow).
\item Write down the relevant equation and identify learned vs fixed terms.
\item Consider numerical implications for training or sampling.
\end{itemize}
\end{block}
\end{frame}

\begin{frame}{Concept check 66: Langevin}
\begin{block}{Question}
Why does Langevin dynamics sample from \(p(x)\) when using the exact score?
\end{block}
\begin{block}{Discussion prompts}
\begin{itemize}
\item Relate to the chapter perspective (variational / score / flow).
\item Write down the relevant equation and identify learned vs fixed terms.
\item Consider numerical implications for training or sampling.
\end{itemize}
\end{block}
\end{frame}

\begin{frame}{Concept check 67: annealing}
\begin{block}{Question}
What is the role of decreasing \(\sigma\) in annealed Langevin dynamics?
\end{block}
\begin{block}{Discussion prompts}
\begin{itemize}
\item Relate to the chapter perspective (variational / score / flow).
\item Write down the relevant equation and identify learned vs fixed terms.
\item Consider numerical implications for training or sampling.
\end{itemize}
\end{block}
\end{frame}

\begin{frame}{Concept check 68: Fokker--Planck}
\begin{block}{Question}
What does the second-order term in Fokker--Planck represent physically?
\end{block}
\begin{block}{Discussion prompts}
\begin{itemize}
\item Relate to the chapter perspective (variational / score / flow).
\item Write down the relevant equation and identify learned vs fixed terms.
\item Consider numerical implications for training or sampling.
\end{itemize}
\end{block}
\end{frame}

\begin{frame}{Concept check 69: reverse drift}
\begin{block}{Question}
Why must the reverse drift include the score term?
\end{block}
\begin{block}{Discussion prompts}
\begin{itemize}
\item Relate to the chapter perspective (variational / score / flow).
\item Write down the relevant equation and identify learned vs fixed terms.
\item Consider numerical implications for training or sampling.
\end{itemize}
\end{block}
\end{frame}

\begin{frame}{Concept check 70: PF-ODE}
\begin{block}{Question}
Why can an ODE share marginals with an SDE?
\end{block}
\begin{block}{Discussion prompts}
\begin{itemize}
\item Relate to the chapter perspective (variational / score / flow).
\item Write down the relevant equation and identify learned vs fixed terms.
\item Consider numerical implications for training or sampling.
\end{itemize}
\end{block}
\end{frame}

\begin{frame}{Concept check 71: divergence}
\begin{block}{Question}
In CNFs, why does divergence appear in the log-density dynamics?
\end{block}
\begin{block}{Discussion prompts}
\begin{itemize}
\item Relate to the chapter perspective (variational / score / flow).
\item Write down the relevant equation and identify learned vs fixed terms.
\item Consider numerical implications for training or sampling.
\end{itemize}
\end{block}
\end{frame}

\begin{frame}{Concept check 72: flow matching}
\begin{block}{Question}
What does flow matching regress and why can it avoid likelihood evaluation?
\end{block}
\begin{block}{Discussion prompts}
\begin{itemize}
\item Relate to the chapter perspective (variational / score / flow).
\item Write down the relevant equation and identify learned vs fixed terms.
\item Consider numerical implications for training or sampling.
\end{itemize}
\end{block}
\end{frame}

\begin{frame}{Concept check 73: parameterizations}
\begin{block}{Question}
Give one practical advantage of \(\epsilon\)-prediction vs \(x_0\)-prediction.
\end{block}
\begin{block}{Discussion prompts}
\begin{itemize}
\item Relate to the chapter perspective (variational / score / flow).
\item Write down the relevant equation and identify learned vs fixed terms.
\item Consider numerical implications for training or sampling.
\end{itemize}
\end{block}
\end{frame}

\begin{frame}{Concept check 74: discretization}
\begin{block}{Question}
How does step size affect sampling quality and computational cost?
\end{block}
\begin{block}{Discussion prompts}
\begin{itemize}
\item Relate to the chapter perspective (variational / score / flow).
\item Write down the relevant equation and identify learned vs fixed terms.
\item Consider numerical implications for training or sampling.
\end{itemize}
\end{block}
\end{frame}

\begin{frame}{Concept check 75: conditioning}
\begin{block}{Question}
Why is time conditioning essential for a single network across all noise levels?
\end{block}
\begin{block}{Discussion prompts}
\begin{itemize}
\item Relate to the chapter perspective (variational / score / flow).
\item Write down the relevant equation and identify learned vs fixed terms.
\item Consider numerical implications for training or sampling.
\end{itemize}
\end{block}
\end{frame}

\begin{frame}{Concept check 76: KL directions}
\begin{block}{Question}
Why does \(\mathrm{KL}(p_{\text{data}}\|p_\theta)\) penalize missing modes more than \(\mathrm{KL}(p_\theta\|p_{\text{data}})\)?
\end{block}
\begin{block}{Discussion prompts}
\begin{itemize}
\item Relate to the chapter perspective (variational / score / flow).
\item Write down the relevant equation and identify learned vs fixed terms.
\item Consider numerical implications for training or sampling.
\end{itemize}
\end{block}
\end{frame}

\begin{frame}{Concept check 77: Markov property}
\begin{block}{Question}
In DDPM forward process, why is \(q(x_t\mid x_{t-1})\) chosen Markovian?
\end{block}
\begin{block}{Discussion prompts}
\begin{itemize}
\item Relate to the chapter perspective (variational / score / flow).
\item Write down the relevant equation and identify learned vs fixed terms.
\item Consider numerical implications for training or sampling.
\end{itemize}
\end{block}
\end{frame}

\begin{frame}{Concept check 78: Gaussianity}
\begin{block}{Question}
Which assumptions make \(q(x_t\mid x_0)\) Gaussian?
\end{block}
\begin{block}{Discussion prompts}
\begin{itemize}
\item Relate to the chapter perspective (variational / score / flow).
\item Write down the relevant equation and identify learned vs fixed terms.
\item Consider numerical implications for training or sampling.
\end{itemize}
\end{block}
\end{frame}

\begin{frame}{Concept check 79: regression view}
\begin{block}{Question}
Why does the ELBO reduce to an MSE loss under \(\epsilon\)-prediction?
\end{block}
\begin{block}{Discussion prompts}
\begin{itemize}
\item Relate to the chapter perspective (variational / score / flow).
\item Write down the relevant equation and identify learned vs fixed terms.
\item Consider numerical implications for training or sampling.
\end{itemize}
\end{block}
\end{frame}

\begin{frame}{Concept check 80: score meaning}
\begin{block}{Question}
Interpret \(s(x)=\nabla_x\log p(x)\) geometrically.
\end{block}
\begin{block}{Discussion prompts}
\begin{itemize}
\item Relate to the chapter perspective (variational / score / flow).
\item Write down the relevant equation and identify learned vs fixed terms.
\item Consider numerical implications for training or sampling.
\end{itemize}
\end{block}
\end{frame}

\begin{frame}{Project ideas}
\begin{itemize}
\item Toy DDPM on 2D mixtures; visualize trajectories and learned denoiser.
\item Train a time-conditioned score model; sample via reverse SDE and PF-ODE.
\item Compare solvers: Euler vs higher-order on PF-ODE; measure FID-like metrics on toy data.
\item Flow matching vs DDPM on toy transport; compare sample quality and speed.
\end{itemize}
\end{frame}

\begin{frame}{Wrap-up}
\begin{itemize}
\item Chapters 1--6 establish diffusion models as \emph{probability transport}.
\item Three perspectives agree: learn a time-dependent vector field.
\item Next topics (later chapters): conditioning/guidance, inverse problems, acceleration, optimal transport links.
\end{itemize}
\centering{\Large Questions?}
\end{frame}

\end{document}
